\documentclass[10pt, aspectratio=169]{beamer}
\input{../../_en_preambulo_beamer}
% Comandos y macros estadísticas compartidas para las presentaciones Beamer en inglés (Metropolis)

% Conjuntos numéricos básicos
\newcommand{\R}{\mathbb{R}}
\newcommand{\N}{\mathbb{N}}
\newcommand{\Q}{\mathbb{Q}}
\newcommand{\Z}{\mathbb{Z}}
\newcommand{\set}[1]{\left\{ #1 \right\}}
\newcommand{\abs}[1]{\left|#1\right|}
\newcommand{\norm}[1]{\left\| #1 \right\|}

% Comandos estadísticos y probabilísticos del curso
\newcommand{\Var}{\operatorname{Var}}
\newcommand{\cov}{\operatorname{Cov}}
\newcommand{\corr}{\rho}
\newcommand{\comb}[2]{\begin{pmatrix} #1 \\ #2 \end{pmatrix}}
\newcommand{\s}{\sigma}
\newcommand{\particion}{\mathfrak{P}}
\newcommand{\card}[1]{\#\left( #1 \right)}
\newcommand{\seq}[3]{\set{#1}_{#2}^{#3}}
\newcommand{\E}{\mathbb{E}}
\newcommand{\Prob}{\mathbb{P}}

% Entornos pedagógicos y cajas en inglés académico natural (soportando tanto nombres en inglés como en español)
\newenvironment{discussion}{%
  \begin{alertblock}{Discussion Question}%
}{%
  \end{alertblock}%
}
\newenvironment{discusion}{%
  \begin{alertblock}{Discussion Question}%
}{%
  \end{alertblock}%
}

\newenvironment{reflection}{%
  \begin{alertblock}{Classroom Reflection}%
}{%
  \end{alertblock}%
}
\newenvironment{reflexion}{%
  \begin{alertblock}{Classroom Reflection}%
}{%
  \end{alertblock}%
}

\newenvironment{paradox}{%
  \begin{alertblock}{Exploratory Paradox}%
}{%
  \end{alertblock}%
}
\newenvironment{paradoja}{%
  \begin{alertblock}{Exploratory Paradox}%
}{%
  \end{alertblock}%
}

\newenvironment{motivation}{%
  \begin{exampleblock}{Motivation: Data Science in Practice}%
}{%
  \end{exampleblock}%
}
\newenvironment{motivacion}{%
  \begin{exampleblock}{Motivation: Data Science in Practice}%
}{%
  \end{exampleblock}%
}

\newenvironment{quickchallenge}{%
  \begin{block}{Quick Team Challenge}%
}{%
  \end{block}%
}
\newenvironment{retoclase}{%
  \begin{block}{Quick Team Challenge}%
}{%
  \end{block}%
}


\title[$Z$/$t$ Tests for Means]{Section 07.02: $Z$ and $t$ Tests for One- and Two-Sample Means}
\subtitle{Independent, Homoscedastic, Welch, and Paired Designs}
\author[J. Castillo Colmenares]{Juliho Castillo Colmenares}
\institute{Tecnológico de Monterrey}
\date{\vspace{-1.5cm}}

\begin{document}

% Slide 1: Title
\begin{frame}[plain]
  \titlepage
\end{frame}

% Slide 2: Roadmap
\begin{frame}{Roadmap --- Chapter 07: Hypothesis Testing}
  \footnotesize
  Where are we in the modular development of the chapter? \pause
  \begin{itemize}\itemsep=0.08em
    \item \textbf{07.01} Foundations: $H_0$ vs. $H_1$, Type I and II Errors, and Power \pause
    \item \textbf{07.02} $Z$ and $t$ Tests for One- and Two-Sample Means \emph{(Today)} \pause
    \item \textbf{07.03} Chi-Squared Goodness-of-Fit Tests \pause
    \item \textbf{07.04} Contingency Tables and Independence Tests
  \end{itemize}
  \vspace{-0.05cm}
  \begin{block}{Session Objective}
    Extend the logical framework of Section 07.01 (hypothesis formulation, $\alpha$, $\beta$, power) to the concrete tests for one and two means: the one-sample $t$-test, the homoscedastic two-independent-sample case (pooled variance), the heteroscedastic case (Welch), and the paired-sample case.
  \end{block}
\end{frame}

% Slide 3: Motivation
\begin{frame}{From One Mean to Comparing Groups}
  \footnotesize
  \begin{motivacion}
    In Section 07.01 we applied the hypothesis-testing framework to a single mean with known $\sigma$. In practice, we almost never know $\sigma$, and the real goal is frequently to \emph{compare} two groups: is algorithm A faster than B? Does a treatment improve the performance of the same subjects?
  \end{motivacion}

  \vspace{0.15cm}
  \pause
  \begin{itemize}\itemsep=0.1cm
    \item \textbf{Unknown $\sigma$ (one sample):} replace $\sigma$ with $S$ and use the $t$-distribution with $n-1$ df. \pause
    \item \textbf{Two independent samples:} the statistic depends on whether population variances are equal or different. \pause
    \item \textbf{Paired samples:} algebraically reduce to a one-sample $t$-test on the differences.
  \end{itemize}
\end{frame}

% Slide 4: One-sample t-test
\begin{frame}{One-Sample $t$-Test (Unknown $\sigma$)}
  \footnotesize
  \begin{block}{Test Statistic}
    $$t = \frac{\bar X - \mu_0}{S/\sqrt n} \ \sim\ t_{n-1} \quad \text{under } H_0: \mu=\mu_0$$
  \end{block}
  \pause

  \vspace{0.1cm}
  \begin{alertblock}{Comparison with Section 07.01}
    The structure is identical to the $Z$-test, except that: (1) $\sigma$ is replaced by its sample estimator $S$, and (2) the reference distribution moves from $N(0,1)$ to $t_{n-1}$, which is more spread out to compensate for the additional uncertainty of estimating $\sigma$.
  \end{alertblock}
\end{frame}

% Slide 5: Two independent means (pooled and Welch)
\begin{frame}{Two Independent Means: Pooled Variance vs. Welch}
  \footnotesize
  \begin{block}{Homoscedastic Case ($\sigma_1^2=\sigma_2^2$): Pooled Variance}
    $$t = \frac{(\bar X_1-\bar X_2)-\Delta_0}{S_p\sqrt{1/n_1+1/n_2}}, \qquad S_p^2 = \frac{(n_1-1)S_1^2+(n_2-1)S_2^2}{n_1+n_2-2}, \qquad \nu=n_1+n_2-2$$
  \end{block}
  \pause

  \vspace{0.1cm}
  \begin{alertblock}{Heteroscedastic Case ($\sigma_1^2\neq\sigma_2^2$): Welch-Satterthwaite}
    $$t = \frac{(\bar X_1-\bar X_2)-\Delta_0}{\sqrt{S_1^2/n_1+S_2^2/n_2}}, \qquad \nu_{\text{Welch}} = \frac{(S_1^2/n_1+S_2^2/n_2)^2}{\frac{(S_1^2/n_1)^2}{n_1-1}+\frac{(S_2^2/n_2)^2}{n_2-1}}$$
  \end{alertblock}
\end{frame}

% Slide 6: Paired samples
\begin{frame}{Paired Samples: Reduction to a Single Sample}
  \footnotesize
  \begin{block}{Difference Variable}
    For repeated measurements on the same subjects ($X_{1i}, X_{2i}$), we define $D_i=X_{1i}-X_{2i}$ and reduce the contrast on $\mu_1-\mu_2$ to a one-sample $t$-test on $\mu_D$:
    $$t = \frac{\bar D - \Delta_0}{S_D/\sqrt n} \ \sim\ t_{n-1}$$
  \end{block}
  \pause

  \vspace{0.1cm}
  \begin{alertblock}{Why Do We Pair?}
    Pairing removes \emph{between-subject} variability (each subject is its own control), drastically increasing the test's power by shrinking the standard error relative to an independent-samples design.
  \end{alertblock}
\end{frame}

% Slide 7: Decision tree
\begin{frame}{Which Test to Choose? A Decision Tree}
  \footnotesize
  \begin{itemize}\itemsep=0.12cm
    \item \textbf{One sample only?} Use the $t$-test (or $Z$-test if $\sigma$ is known) from Section 07.01. \pause
    \item \textbf{Two samples, same subjects measured twice?} Use the paired $t$-test on $D_i$. \pause
    \item \textbf{Two independent samples, similar variances?} Use the pooled-variance $t$-test. \pause
    \item \textbf{Two independent samples, clearly different variances?} Use Welch's $t$-test.
  \end{itemize}
\end{frame}

% Slide 8: Python Lab (Block 1)
\begin{frame}[fragile]{Python Lab: One-Sample $t$-Test}
  \scriptsize
  $t$-test with unknown $\sigma$, verified against \texttt{scipy.stats.ttest\_1samp} (\texttt{07.02\_z\_t\_tests\_means.py}):
  {\lstset{basicstyle=\fontsize{4pt}{4.8pt}\selectfont\ttfamily,aboveskip=0.1em,belowskip=0.1em}
  \lstinputlisting[firstline=17, lastline=37]{../../code/07_pruebas_hipotesis/07.02_z_t_tests_means.py}}
\end{frame}

% Slide 9: Python Lab (Block 2)
\begin{frame}[fragile]{Python Lab: Pooled Variance vs. Welch}
  \scriptsize
  Direct comparison of both statistics under homoscedasticity and heteroscedasticity:
  {\lstset{basicstyle=\fontsize{3.6pt}{4.3pt}\selectfont\ttfamily,aboveskip=0.05em,belowskip=0.05em}
  \lstinputlisting[firstline=40, lastline=77]{../../code/07_pruebas_hipotesis/07.02_z_t_tests_means.py}}
\end{frame}

% Slide 10: Python Lab (Block 3)
\begin{frame}[fragile]{Python Lab: Paired $t$-Test}
  \scriptsize
  Reducing a paired design to a one-sample $t$-test on the differences:
  {\lstset{basicstyle=\fontsize{4pt}{4.8pt}\selectfont\ttfamily,aboveskip=0.1em,belowskip=0.1em}
  \lstinputlisting[firstline=80, lastline=101]{../../code/07_pruebas_hipotesis/07.02_z_t_tests_means.py}}
\end{frame}

% Slide 11: Terminal Output
\begin{frame}[fragile]{Python Lab: Terminal Output}
  \scriptsize
  Standard output produced when running the Python lab script:
  \vspace{0.02cm}
  \begin{block}{Standard Console Output}
    \fontsize{4pt}{5pt}\selectfont
    \begin{verbatim}
=== Block 1: One-Sample t-Test ===
t statistic = 1.7050, df=48, two-tailed p-value = 0.0947

=== Block 2: Pooled vs. Welch t-Test ===
Homoscedastic: Pooled t=-1.9815, df=32, p=0.0562
Heteroscedastic: Welch t=4.2026, df=16.60, p=0.0006

=== Block 3: Paired t-Test ===
t statistic = 6.3743, df=11, one-tailed p-value = 0.000026
    \end{verbatim}
  \end{block}
\end{frame}


% Slide 16: Closing
\begin{frame}{Section Closing: $Z$/$t$ Tests for Means}
  \begin{reflexion}
    Comparing means first requires diagnosing the structure of the problem --- one sample or two? independent or paired? equal or unequal variances? --- before mechanically applying a formula. Choosing the wrong statistic (e.g., pooled variance when variances are clearly different) can invalidate the inference.
  \end{reflexion}

  \vspace{0.2cm}
  \pause
  \begin{block}{Key Takeaways}
    \begin{itemize}\itemsep=0.08cm
      \item \textbf{One sample, unknown $\sigma$:} $t=(\bar X-\mu_0)/(S/\sqrt n)$, $\nu=n-1$.
      \item \textbf{Two independent samples:} pooled variance if $\sigma_1^2=\sigma_2^2$; Welch otherwise.
      \item \textbf{Paired samples:} reduce to a one-sample $t$-test on the differences $D_i$.
    \end{itemize}
  \end{block}
\end{frame}

% Slide 17: Modular Perspective
\begin{frame}{Modular Perspective --- The Next Challenge}
  \scriptsize
  \begin{retoclase}
    So far we have tested hypotheses about means of \emph{continuous} variables. Section 07.03 shifts terrain: we will test whether complete \emph{categorical} data fits a proposed theoretical distribution, using the $\chi^2$ goodness-of-fit test --- the fundamental tool for validating complete discrete models, not just a single isolated parameter.
  \end{retoclase}

  \vspace{0.08cm}
  \pause
  \begin{alertblock}{Recommended Preparation}
    \begin{itemize}\itemsep=0.06cm
      \item Review the $\chi^2$ distribution and its degrees of freedom (Section 05.03).
      \item Reflect on why estimating parameters from the sample itself reduces the available degrees of freedom.
    \end{itemize}
  \end{alertblock}
\end{frame}

\end{document}
